\section{Citation Verification over a Legal Database}
\label{sec:verify}

The verifier answers a deceptively simple question: \emph{does article
$\langle$num$\rangle$ of $\langle$code$\rangle$ exist, and is it in force?}
Answering it soundly \emph{and} at interactive latency over Canutes is the
technical core of the system.

\subsection{Finding the needle in 253 tables}
Because Canutes ships without documented foreign keys and with opaque,
historically accreted table names, \emph{locating} the authoritative article
table is itself work. We recover a usable map of the schema
\emph{offline}---entity-and-relationship inference over the 253 tables,
cross-checked against L\'egifrance's identifier conventions
(\texttt{LEGITEXT}/\texttt{LEGIARTI})---and pin verification to
\texttt{legifrance.article}, whose rows are the versioned records of
\S\ref{sec:background}. Paying this cost once, offline, keeps the online path a
single, well-understood query rather than an exploratory crawl across an
unknown schema.

\subsection{Extraction and normalization}
Citations are first lifted from the draft by a tolerant matcher keyed on the
French legal idiom (\emph{``article~\ldots\ du Code~\ldots''}), capturing the
article number and the code name. L\'egifrance, however, encodes article
numbers \emph{inconsistently}: the same article surfaces as \texttt{L3121-27},
\texttt{L.\,3121-27}, or \texttt{L. 3121-27}, with or without a space after the
book letter and with or without the period. A naive existence predicate that
tries to absorb this variation with a pattern
match---\texttt{num \textasciitilde{} '.*3121-27.*'}---forces a sequential scan
and a regular-expression evaluation over the \emph{entire}
\texttt{legifrance.article} table; in our environment this reliably hit the
statement timeout.

Instead, we push the variation to the \emph{client}: a citation is expanded
into a small, closed set of \emph{exact} string variants
(Algorithm~\ref{alg:variants}), which are then matched with \emph{equality}
against the indexable \texttt{num} column. Equality against a bounded set turns
a full-table regex scan into an index probe.

\begin{algorithm}[t]
\caption{Variant expansion: turn one written citation number into a closed set
of exact strings for index-friendly equality matching.}
\label{alg:variants}
\begin{algorithmic}[1]
\Function{Variants}{\textit{num}}
  \State \textit{core} $\gets$ strip spaces and periods from \textit{num}
         \Comment{e.g.\ \texttt{L3121-27}}
  \State \Return the closed set re-inserting the optional period and space
         after the leading letter(s), e.g.\ $\{$\texttt{L3121-27},
         \texttt{L.3121-27}, \texttt{L. 3121-27}, \texttt{L 3121-27}$\}$
\EndFunction
\end{algorithmic}
\end{algorithm}

\begin{lstlisting}[style=lr, language=SQL]
-- variants = {"L. 3121-27", "L.3121-27", "L3121-27", ...}
SELECT id, num, data
FROM   legifrance.article
WHERE  num = ANY(%s)          -- equality, index-friendly
  AND  data::text ILIKE %s    -- code name appears in the article JSON
LIMIT  200;
\end{lstlisting}

\subsection{Code disambiguation and accents}
The code name (e.g., \emph{Code du travail}) is first filtered in SQL by
\texttt{ILIKE} against the article's JSON payload---cheap, because the variant
equality has already reduced the candidate set to a handful of rows. It is then
re-checked in Python under an accent-insensitive normalization, because
PostgreSQL \texttt{ILIKE} folds case but \emph{not} diacritics: without this
second stage, \emph{``Code de la s\'ecurit\'e sociale''} and \emph{``Code de la
securite sociale''} would be treated as different codes. This two-stage
filter---coarse in SQL, exact in the client---keeps the query cheap while
remaining correct on accented input.

\subsection{State-aware selection and provenance}
Among the surviving rows the verifier prefers the \texttt{VIGUEUR} (in-force)
version. When one is found it returns the canonical \texttt{LEGIARTI}
identifier, from which a real, clickable L\'egifrance URL is constructed, plus
a short textual excerpt for display. The URL and excerpt are the
\emph{provenance} that turns a released citation into something a user can
independently check; they are never synthesized, only copied from the row. If
no surviving row matches the requested code, the verifier returns
\texttt{exists=false}, which the pipeline treats as a refusal trigger.
Crucially, \emph{every} failure mode---no row, no in-force row, timeout,
malformed number---maps to ``not verified''. That total mapping is what makes
the layer fail-closed: there is no code path in which uncertainty is reported
as success.

\subsection{Cost and offline determinism}
Because matching is an equality probe against a bounded variant set followed by
an \texttt{ILIKE} over a handful of candidates, per-citation verification is a
small, bounded query rather than a corpus scan, and the checks for the several
citations in one answer are independent and can be issued concurrently. For
demonstrations and continuous integration without network access, a mock
verifier answers from a small local set of articles known to be in force,
exercising the \emph{same} code path and decision logic as the live backend.
This yields a deterministic, reproducible refusal for a fabricated
article---exactly what the executable specification of \S\ref{sec:bdd} needs to
turn trust into a repeatable test.

\subsection{Verifier backends, one contract}
\label{sec:verify-backends}
The verifier is an \emph{interface} with interchangeable implementations that
all honour the same fail-closed contract: (i)~a \emph{mock} over a small
in-force set, for offline determinism; (ii)~a \emph{direct} PostgreSQL client
issuing the query above; (iii)~an \emph{MCP} client that calls a tool server
(``Moulineuse'') exposing SQL/JS over Canutes; and (iv)~a read-only
\emph{PostgREST} fallback for the public tables. Whichever backend is
configured, an unresolved citation or \emph{any} error maps to ``not
verified''. The guarantee is thus a property of the contract, not of a
particular data path---so the trust layer survives changes in how the corpus is
reached. Putting the pieces together, one verification runs as:

\begin{lstlisting}[style=lr]
verify("L. 3121-27", "Code du travail")
  variants -> {L3121-27, L.3121-27, L. 3121-27, ...}
  SELECT id,num,data FROM legifrance.article
    WHERE num = ANY(variants)
      AND data::text ILIKE '%travail%';
  rows     -> [ LEGIARTI...815 (VIGUEUR), ... ]
  select   -> VIGUEUR  =>  exists = true
  url      -> legifrance.gouv.fr/.../LEGIARTI...815
  excerpt  -> "La duree legale du travail effectif..."
\end{lstlisting}
